\documentclass[11pt]{letter}

\usepackage[margin=1in]{geometry}
\usepackage{hyperref}
\usepackage{parskip} % no indentation, space between paragraphs

\signature{Vedaang Chopra}
\address{
Vedaang Chopra\\
Atlanta, Georgia, United States\\
\href{mailto:vedaangchopra@gatech.edu}{vedaangchopra@gatech.edu} \\
\href{https://www.linkedin.com/in/vedaang-chopra/}{linkedin.com/in/vedaang-chopra}
}

\begin{document}

\begin{letter}{Hiring Committee\\
The Allen Institute for Artificial Intelligence (Ai2)\\
Seattle, WA}

\opening{Dear Hiring Committee,}

I am currently a Master’s student in Computer Science at the Georgia Institute of Technology, where my focus is on AI research spanning large language models, multimodal learning, and agentic systems. Prior to graduate school, I spent nearly five years in industry working on applied AI systems at Fortinet, which gave me practical experience building, evaluating, and deploying machine learning models at scale. I am applying for the Summer 2026 Research Internship at the Allen Institute for Artificial Intelligence.

I was first introduced to Ai2’s work through my graduate coursework in Vision--Language Models, where I presented recent Ai2 research as part of a seminar, including work on Molmo. Studying this work in depth left a strong impression on me, particularly Ai2’s emphasis on open research practices and the release of high-quality public resources such as PixMo. I greatly value Ai2’s commitment to openness, reproducibility, and building artifacts that enable the broader research community.

Inspired by these ideas, I have been developing a research project called Edge-Glass, which explores efficient and modular vision--language systems. In this work, I combine a standard mid-scale language model with a lightweight recursive decoding approach and a CLIP-based vision encoder post-trained on PixMo-style data. The project focuses on understanding trade-offs between model capacity, efficiency, and generalization, and includes a simple routing mechanism to study adaptive model selection. This work draws inspiration from Ai2’s broader research direction in modular model design, evaluation, and agentic capabilities.

My research interests broadly lie in foundational AI research, including model pre-training, post-training, and the development of end-to-end agentic systems that can reason, use tools, and operate robustly in complex environments. I am particularly interested in research that combines system-level thinking with careful empirical evaluation and that contributes reusable insights, benchmarks, or open artifacts. I enjoy research-driven work and am eager to learn from and contribute to Ai2’s collaborative research environment.

I would be excited for the opportunity to contribute to Ai2’s ongoing research efforts while learning from its researchers and engineers. Thank you for your time and consideration.

\closing{Sincerely,}

\end{letter}
\end{document}
